\subsection{The interfaces completed in the sparse restoration}

The preceding propositions prove more than spectral consistency:
the weighted projection has bounded tilted mean and variance, and the
sharp tail permits exceptional coefficients whose expected hit mass is
small. Three additional estimates turn those tools into the actual
algorithm theorem. First, \cref{restore:tail-energy-lemma} proves the
needed truncated energy on the full and recomputed deleted graphs.
Second, \cref{restore:loo-scalar,restore:loo-corrected} compare the
actual full-graph score to the independent deleted score, including
the empirical floor, exact degree correction, and hollowing, at
$O(1)$ cost using row bounds only on incident coordinates.
Third, \cref{restore:fit} controls initialization and fitted likelihood
contrasts by aggregate weighted errors, while retaining exceptional
vertices. The invariant-basin argument covers finite hard refits
without a union over their number.

For $d=o(\log n)$ the already proved $n^{-r}$ geometry failure is
negligible relative to $d^{-1/2}e^{-C_Id}$ for any fixed $r>0$.
This explains why the regularized norm event needs no stronger
probability in that range. Nor should one declare failure of the entire
graph whenever any vertex has an unusually small degree. The
truncated-score comparison shows how local exceptional coordinates can
instead be paid when the tested row encounters them.

Zhang's spectral clustering proof
\cite[Theorem 3 and Corollary 1, p.~7326]{zhang} establishes exponentially
small expected truncated row energy and exceptional-row counts for
the trimmed adjacency eigenspace. That result provides a concrete
model for \eqref{eq:sparse-row-count-target}; its eigenspace and
normalization differ from the floor-normalized DC construction here.
Its final leading-exponent conclusion does not, by itself, imply
the bounded perturbation cost and sharp prefactor required here.
The intended nontrivial sparse target on the present fixed-separation,
bounded-theta class is $d=np\to\infty$, together with the stated
Bernoulli remainder condition. A vanishing signal gap or degree
parameters approaching zero requires additional quantitative work.
The completed proof in \cref{sec:sparse-restored} establishes each
of these estimates. It does not replace them by an assumption that
the desired constant risk ratio already holds.

